% Additive proof fragment.  All matrices use literal powers of S.
% Provider bibliography key: prog:arithmetic.
\section{The full arithmetic action boundary in the original quotient}

The polynomial quotient and positive arithmetic moments in this section are
the original ones of the arithmetic construction
\cite{prog:arithmetic}.  The calculation below uses the full relation ideal,
the original physical line, and the attained quotient metric.  Every identity
needed for the boundary calculation is proved here.

\subsection{Spaces, moments, and the attained section}

In the actual packet \(a=k+4\), with \(k\) in its original guarded domain;
retain that domain and its original \(0<\delta<1/2\) and \(\gamma>2\).
The finite identities below in fact use only that \(a\) is a positive
integer and \(\delta,\gamma>0\).  Put
\begin{equation}
 \beta_{rt}=\frac a2+(2r-a)\delta+i(2t-a)\gamma,\qquad
 \chi(S)=\prod_{r=0}^{a}\prod_{t=0}^{a}(S-\beta_{rt}),\qquad
 q=(a+1)^2 .
 \tag{NAB1}\label{eq:NAB1}
\end{equation}
These roots are pairwise distinct: equality of real parts gives equality
of \(r\), and equality of imaginary parts then gives equality of \(t\).
No root or coefficient of \(\chi\) is discarded.  Write
\(\mathcal P_j=\{p\in\mathbb C[S]:\deg p\leq j\}\), with
\(\mathcal P_{-1}=\{0\}\), and
\(\mathcal E=\mathbb C[S]/(\chi)\).  Coordinates in
\(\mathcal P_j\) are the literal coefficients of
\(1,S,\ldots,S^j\); coordinates in \(\mathcal E\) are the corresponding
classes of \(1,S,\ldots,S^{q-1}\).

Fix an integer \(N\geq q-1\).  Use the original positive measure
\(d\mu_N(y)\) on the physical line \(S=a/2+iy\), with all of its original
components and both integration regions.  The subscript allows the measure
to depend on the fixed endpoint \(N\).  Within this calculation the measure
is held fixed.  Its moment matrix through degree \(N+1\) is
\begin{equation}
 H_j^{[N]}=
 \left(
 \int_{\mathbb R}
    \overline{(a/2+iy)^r}(a/2+iy)^t\,d\mu_N(y)
 \right)_{0\leq r,t\leq j},
 \qquad 0\leq j\leq N+1 .
 \tag{NAB2}\label{eq:NAB2}
\end{equation}
All moments displayed here are finite and \(H_{N+1}^{[N]}\) is positive
definite in the original arithmetic construction.  Equivalently, the
integral of \(|p(a/2+iy)|^2\) is strictly positive for every nonzero
\(p\in\mathcal P_{N+1}\).  The superscript records that
\(H_N^{[N]}\) and \(H_{N+1}^{[N]}\) use the same measure; an endpoint change
of measure is not implicit in this notation.  In the remainder of this
section write these matrices as \(H_N\) and \(H_{N+1}\).

The Hermitian product is conjugate linear in its first entry:
\(\langle f,g\rangle_N=\int\overline{f(a/2+iy)}g(a/2+iy)d\mu_N(y)\).
For \(f,g\in\mathcal P_N\), the identity
\(\overline S+S=a\) on this actual line gives
\begin{equation}
 \langle Sf,g\rangle_N+\langle f,Sg\rangle_N
       =a\langle f,g\rangle_N .
 \tag{NAB3}\label{eq:NAB3}
\end{equation}
The two left-hand products use the moment matrix through degree \(N+1\).
Thus (NAB3) retains the physical coordinate rather than replacing the
action by an abstract self-adjoint multiplication operator.

Let \(\rho_N:\mathcal P_N\longrightarrow\mathcal E\) be reduction modulo
\(\chi\).  Polynomial division by the original monic \(\chi\) proves
\begin{equation}
 \ker\rho_N=\mathcal R_N:=\chi\mathcal P_{N-q},\qquad
 \mathcal P_N/\mathcal R_N \xrightarrow[\rho_N]{\ \cong\ }\mathcal E .
 \tag{NAB4}\label{eq:NAB4}
\end{equation}
For \(N=q-1\), \(\mathcal R_N=\{0\}\).
Let \(E_N:\mathbb C^q\to\mathbb C^{N+1}\) be the coefficient inclusion of
\(\mathcal P_{q-1}\) into \(\mathcal P_N\).  If \(N\geq q\), let \(W_N\)
have columns the coefficient vectors of
\(\chi,S\chi,\ldots,S^{N-q}\chi\).  Its columns are linearly independent
because multiplication by a nonzero polynomial is injective.
Consequently \(W_N^*H_NW_N\) is positive definite.  The attained section is
\begin{equation}
 T_N=E_N-W_N(W_N^*H_NW_N)^{-1}W_N^*H_NE_N,\qquad
 G_N=T_N^*H_NT_N ;
 \quad T_{q-1}=E_{q-1}.
 \tag{NAB5}\label{eq:NAB5}
\end{equation}
Here \(T_N:\mathcal E\to\mathcal P_N\) uses the fixed coefficient
coordinates, and \(G_N\) is its \(q\)-by-\(q\) attained metric.
Indeed, \(\rho_NT_N=I_{\mathcal E}\) by (NAB4), and direct multiplication
gives \(W_N^*H_NT_N=0\).  Every representative of \(v\in\mathcal E\)
is \(T_Nv+W_Nz\) for a unique \(z\).  Orthogonality therefore proves
\[
 \|T_Nv+W_Nz\|_N^2
 =v^*G_Nv+z^*W_N^*H_NW_Nz.
\]
This proves attainment, uniqueness, and the original minimum property
without a change of relation space.  It also proves \(G_N>0\), since
\(T_Nv\ne0\) whenever \(v\ne0\).

\subsection{The exact one-column boundary and its Hermitian form}

Let \(A:\mathcal E\to\mathcal E\) be multiplication by the original \(S\)
modulo \(\chi\).  Write
\(\chi(S)=S^q+\sum_{j=0}^{q-1}c_jS^j\).
In the fixed ordered basis,
\begin{equation}
 A e_j=e_{j+1}\quad(0\leq j<q-1),\qquad
 A e_{q-1}=-\sum_{j=0}^{q-1}c_je_j .
 \tag{NAB6}\label{eq:NAB6}
\end{equation}
Embed \(T_N\) into \(\mathcal P_{N+1}\) by appending a zero top
coefficient, writing this embedding as \(\widetilde T_N\).
Define the row, polynomial column, and vector
\begin{equation}
 \ell_N=[S^N]T_N:\mathcal E\longrightarrow\mathbb C,\qquad
 b_N=\operatorname{coeff}_{0,\ldots,N+1}
             \bigl(\chi(S)S^{N+1-q}\bigr),\qquad
 \xi_N=\widetilde T_N^*H_{N+1}b_N\in\mathbb C^q .
 \tag{NAB7}\label{eq:NAB7}
\end{equation}
The notation \([S^N]\) means the literal coefficient functional, with no
factorial, root evaluation, or rescaling.

Let \(M_N:\mathcal P_N\to\mathcal P_{N+1}\) denote multiplication by \(S\)
in those same coordinates.  There is a linear map
\(Z_N:\mathcal E\to\mathcal R_N\), embedded in \(\mathcal P_{N+1}\),
such that
\begin{equation}
 M_NT_N-\widetilde T_N A=b_N\ell_N+Z_N,\qquad
 \widetilde T_N^*H_{N+1}Z_N=0.
 \tag{NAB8}\label{eq:NAB8}
\end{equation}
To prove the first identity on \(v\in\mathcal E\), reduce
\(S(T_Nv)-T_N(Av)\) modulo \(\chi\).
The result is \(Av-Av=0\), so the polynomial is divisible by \(\chi\).
Its degree is at most \(N+1\), and its coefficient of \(S^{N+1}\) is
exactly \(\ell_Nv\).  Monicity of the full \(\chi\) now shows that
subtracting \(\chi S^{N+1-q}\ell_Nv\) leaves an element of
\(\chi\mathcal P_{N-q}\).
This subtraction is linear in \(v\), giving \(Z_N\).
If \(N=q-1\), the remainder has degree less than \(q\) and is divisible
by \(\chi\), hence is zero.
Otherwise the second identity is exactly the orthogonality in (NAB5):
both factors of this product have degree at most \(N\), so the use of
\(H_{N+1}\) restricts to the same \(H_N\).

Define the full arithmetic action form
\begin{equation}
 D_N=A^*G_N+G_NA-aG_N .
 \tag{NAB9}\label{eq:NAB9}
\end{equation}
It has the exact boundary formula
\begin{equation}
 \boxed{\displaystyle
 D_N=-\xi_N\ell_N-\ell_N^*\xi_N^*.}
 \tag{NAB10}\label{eq:NAB10}
\end{equation}
For a direct proof, (NAB8) and its orthogonality give
\[
 \widetilde T_N^*H_{N+1}M_NT_N
       =G_NA+\xi_N\ell_N .
\]
Taking conjugate transposes gives the other term.
Applying (NAB3) to every pair of columns of \(T_N\) gives
\[
 aG_N=
 \widetilde T_N^*H_{N+1}M_NT_N+
 (M_NT_N)^*H_{N+1}\widetilde T_N
 =G_NA+A^*G_N+\xi_N\ell_N+\ell_N^*\xi_N^*.
\]
Rearranging proves (NAB10), including both minus signs.
In particular \(\operatorname{rank}D_N\leq2\).
No estimate for a trial section or different denominator was used.

\subsection{Exact inertia, trace, and boundary magnitude}

For every root \(\beta\) in (NAB1), define the original Lagrange
polynomial and its class
\begin{equation}
 L_\beta(S)=\frac{\chi(S)}{(S-\beta)\chi'(\beta)}
       \in\mathcal P_{q-1},\qquad
 v_\beta=[L_\beta]\in\mathcal E .
 \tag{NAB11}\label{eq:NAB11}
\end{equation}
The denominator \(\chi'(\beta)\) is nonzero by the distinctness proved
after (NAB1).  At the original roots \(L_\beta(\beta)=1\) and
\(L_\beta(\beta')=0\) for \(\beta'\ne\beta\).
Thus \(v_\beta\ne0\), and the \(q\) such classes form a basis: a linear
relation among them evaluates at each root to its corresponding
coefficient.  The literal polynomial identity
\((S-\beta)L_\beta=\chi/\chi'(\beta)\) proves the exact morphism
\begin{equation}
 Av_\beta=\beta v_\beta,\qquad
 v_{\beta_{rt}}^*D_Nv_{\beta_{rt}}
      =2(2r-a)\delta\,v_{\beta_{rt}}^*G_Nv_{\beta_{rt}} .
 \tag{NAB12}\label{eq:NAB12}
\end{equation}
In particular the expression is strictly negative at \(r=0\) and
strictly positive at \(r=a\).
A Hermitian form of rank at most two with both a positive and a negative
value must have rank exactly two and one sign of each kind:
\begin{equation}
 \operatorname{inertia}(D_N)=(1,1,q-2).
 \tag{NAB13}\label{eq:NAB13}
\end{equation}
For completeness, this elementary inference can be seen by diagonalizing
the Hermitian form using successive orthogonal eigenvectors.  Such a
basis is obtained by maximizing its real Rayleigh quotient on the compact
unit sphere; differentiation in every real and imaginary tangent
direction makes a maximizing vector an eigenvector.  Its orthogonal
complement is invariant because the matrix is Hermitian, so induction
gives the diagonalization.  A positive value requires a positive
diagonal entry and a negative value requires a negative diagonal entry.
The rank bound allows precisely those two nonzero entries.

The full original grid also gives an exact trace cancellation:
\begin{equation}
 \operatorname{tr}(G_N^{-1}D_N)
 =\overline{\operatorname{tr}A}+\operatorname{tr}A-aq
 =2\operatorname{Re}\sum_{r,t=0}^{a}\beta_{rt}-aq=0 .
 \tag{NAB14}\label{eq:NAB14}
\end{equation}
The first equality uses only cyclicity of finite matrix trace.
For the last equality,
\(\sum_{r=0}^a(2r-a)=0=\sum_{t=0}^a(2t-a)\), so the root sum
is exactly \(aq/2\).
The Lagrange basis in (NAB11) proves that this root sum is the trace of
the actual \(A\); no approximate eigenvalue calculation is involved.

Take the positive Hermitian square root of \(G_N\), and define
\begin{equation}
 B_N=G_N^{-1/2}D_NG_N^{-1/2},\qquad
 u_N=G_N^{-1/2}\xi_N,\qquad
 w_N=G_N^{-1/2}\ell_N^*,\qquad
 c_N=w_N^*u_N=\ell_NG_N^{-1}\xi_N .
 \tag{NAB15}\label{eq:NAB15}
\end{equation}
The square root is obtained by the same finite Hermitian
diagonalization, taking positive square roots of the positive entries of
\(G_N\).  Equations (NAB10) and (NAB14) prove
\begin{equation}
 B_N=-u_Nw_N^*-w_Nu_N^*,\qquad
 -2\operatorname{Re}c_N=\operatorname{tr}B_N=0.
 \tag{NAB16}\label{eq:NAB16}
\end{equation}
Congruence by the invertible \(G_N^{-1/2}\) preserves positive and
negative subspaces, hence (NAB13) applies to \(B_N\) too.
In particular \(u_N,w_N\) are linearly independent.
Set \(b=\|w_N\|>0\), \(e_1=w_N/b\), and write
\[
 u_N=(c_N/b)e_1+\eta e_2,\qquad
 \eta=\left(\|u_N\|^2-|c_N|^2/b^2\right)^{1/2}>0,
\]
where \(e_2\) is the resulting unit vector orthogonal to \(e_1\).
On their span the matrix of \(B_N\) is exactly
\[
 \begin{pmatrix}-2\operatorname{Re}c_N&-b\eta\\
                 -b\eta&0\end{pmatrix},
\]
and on its orthogonal complement it is zero.  Since
\(\operatorname{Re}c_N=0\), its nonzero eigenvalues are
\begin{equation}
 \boxed{\displaystyle
 \lambda_\pm(B_N)=\pm\kappa_N,\qquad
 \kappa_N=
 \sqrt{\|u_N\|^2\|w_N\|^2-
                    (\operatorname{Im}(w_N^*u_N))^2}>0 .}
 \tag{NAB17}\label{eq:NAB17}
\end{equation}
The equivalent formula using only the original coefficient matrices is
\begin{equation}
 \kappa_N^2=
  (\xi_N^*G_N^{-1}\xi_N)
  (\ell_NG_N^{-1}\ell_N^*)
  -|\ell_NG_N^{-1}\xi_N|^2
  =\tfrac12\operatorname{tr}\bigl((G_N^{-1}D_N)^2\bigr).
 \tag{NAB18}\label{eq:NAB18}
\end{equation}
The first equality follows from (NAB15)--(NAB17), and the second follows
because \(G_N^{-1}D_N\) is similar to \(B_N\).
The exact grid eigenclasses in (NAB12) furthermore give the finite bound
\begin{equation}
 \kappa_N\geq2a\delta .
 \tag{NAB19}\label{eq:NAB19}
\end{equation}
Indeed, put \(z=G_N^{1/2}v_{\beta_{0t}}\).  It is nonzero and
\(z^*B_Nz/(z^*z)=-2a\delta\).
Every Rayleigh quotient of the diagonal matrix with entries
\(-\kappa_N,\kappa_N,0,\ldots,0\) lies in
\([-\kappa_N,\kappa_N]\); this proves (NAB19).
The choice \(r=a\) gives the matching positive quotient \(2a\delta\).

If the complete grid trace cancellation is not available, the boundary
identity (NAB10) still has its stated sign whenever (NAB2)--(NAB8) hold,
but the two eigenvalues are instead
\[
 -\operatorname{Re}c_N
 \ \pm\sqrt{\|u_N\|^2\|w_N\|^2-(\operatorname{Im}c_N)^2}.
\]
This follows by taking the characteristic polynomial of the displayed
two-by-two matrix.  In the actual full grid the shift vanishes exactly
by (NAB14).  Thus omission of a root subfamily cannot silently retain
the full-grid trace cancellation.

\subsection{An exact original-coordinate sign diagnostic}

The following small example is an exact algebraic diagnostic only.  It
is not an actual packet or period, and is not an evaluation of the
arithmetic source.  It neither replaces nor relaxes any original
large-degree finite guard.
Take \(a=1\), \(\delta=1/4\), \(\gamma=3\), \(q=4\), \(N=4\),
\(S=1/2+iy\), and the standard real Gaussian measure.  The original
root polynomial and relation space are
\begin{equation}
 \begin{aligned}
 \chi(S)&=
 \bigl((S-\tfrac12-\tfrac14)^2+9\bigr)
 \bigl((S-\tfrac12+\tfrac14)^2+9\bigr)\\
 &=S^4-2S^3+\tfrac{155}{8}S^2-\tfrac{147}{8}S+\tfrac{22185}{256},
 \qquad \mathcal R_4=\operatorname{span}\{\chi\}.
 \end{aligned}
 \tag{NAB20}\label{eq:NAB20}
\end{equation}
The roots are exactly \(1/2\pm1/4\pm3i\), all in the right half-plane.
For \(m\geq0\), put
\(\mu_{2m}=(2m-1)!!\), with \((-1)!!=1\), and
\(\mu_{2m+1}=0\).  These are the Gaussian moments: the odd integral
vanishes under \(y\mapsto-y\), and integration by parts gives
\(\mu_{2m}=(2m-1)\mu_{2m-2}\).
The moment entries, entirely in the original \(S\)-coordinates, are
\begin{equation}
 (H_5)_{mn}
   =\sum_{r=0}^{m}\sum_{t=0}^{n}
    \binom mr\binom nt\,2^{-(m+n-r-t)}
            (-\mathrm i)^r\mathrm i^t\mu_{r+t}.
 \tag{NAB21}\label{eq:NAB21}
\end{equation}
Here \(\mathrm i^2=-1\), and \(0\leq m,n\leq5\).
Every nonzero polynomial has positive squared integral because a
nonzero polynomial has finitely many zeros and the Gaussian density is
strictly positive.  Thus the moment and quotient matrices are positive
definite.

Write
\(c=(22185/256,-147/8,155/8,-2,1)^{\mathsf T}\), and put
\(d=316481153\).
The finite sums (NAB21) give \(c^*H_4c=d/65536\).
Substitution into (NAB5) and (NAB7) gives, exactly,
\begin{equation}
 \begin{aligned}
 \ell_4&=\frac1d(-4407552,-2203776,1749184,3725664),\\
 \xi_4&=\frac1{1024}(0,44548,44548,-51849)^{\mathsf T}.
 \end{aligned}
 \tag{NAB22}\label{eq:NAB22}
\end{equation}
and hence
\begin{equation}
 D_4=-\xi_4\ell_4-\ell_4^{\mathsf T}\xi_4^{\mathsf T}.
 \tag{NAB23}\label{eq:NAB23}
\end{equation}
To specify the remaining finite computation without omitting its
inputs, take \(E_4\) to be the first four columns of the five-by-five
identity, \(T_4=E_4-(65536/d)c(c^*H_4E_4)\), and
\(G_4=T_4^*H_4T_4\), with \(H_4\) the leading block of (NAB21).
Direct substitution and ordinary exact matrix multiplication give
\begin{equation}
 \begin{aligned}
 \operatorname{tr}(G_4^{-1}D_4)&=0,\\
 \frac{\det(zG_4-D_4)}{\det G_4}
      &=z^2\left(z^2-\frac{46738154483830145}{497781812232192}\right),\\
 \kappa_4^2&=\frac{46738154483830145}{497781812232192}.
 \end{aligned}
 \tag{NAB24}\label{eq:NAB24}
\end{equation}
All the matrices and exact rational substitutions for this diagnostic
are recorded by the accompanying
\texttt{exact\_boundary\_example.py} and
\texttt{EXACT\_EXAMPLE.json}.  The general proof is (NAB3)--(NAB19);
the example supplies an independent check of its signs and its use of
the nontrivial attained relation space.

\paragraph{Scope of the calculation.}
Equations (NAB10), (NAB13), and (NAB17)--(NAB19) determine the full
action form and its full-space inertia.  For a specified map
\(J:\mathbb C^d\to\mathcal E\), their exact receiver is
\[
 J^*D_NJ=-(J^*\xi_N)(\ell_NJ)
                 -(\ell_NJ)^*(\xi_N^*J).
\]
This formula carries the original class-position data into the
restricted form.  Evaluating the sign on an individual projected class
requires its actual \(J\), or the corresponding attained projection,
inside that formula.  No individual projected sign, absolute
allocation, or value of the responses \(U,V\) is asserted here.
